\chapter{Contrato operativo}

\section{Función del contrato}
\begin{observation}
El contrato operativo del Tomo II fija las condiciones mínimas bajo las cuales una formulación puede ser tratada como pieza disciplinada del canon intermedio.
\end{observation}

\begin{observation}
Su función es evitar mezcla entre:
\begin{enumerate}
  \item núcleo axiomático,
  \item desarrollo operativo,
  \item laboratorio,
  \item demo,
  \item traducción metodológica a claim.
\end{enumerate}
\end{observation}

\section{Condición de entrada}
\begin{definition}
Una formulación puede entrar al Tomo II si cumple simultáneamente:
\begin{enumerate}
  \item dependencia explícita del núcleo basal integrado en Tomo I;
  \item incremento real de capacidad operativa o jerárquica;
  \item alcance declarado;
  \item compatibilidad con el dominio verificable;
  \item no duplicación innecesaria de primitivas ya fijadas.
\end{enumerate}
\end{definition}

\begin{observation}
La entrada al Tomo II no exige todavía validación experimental completa, pero sí exige disciplina formal suficiente para evitar inflación conceptual.
\end{observation}

\section{Condición de exclusión}
\begin{definition}
No debe entrar al Tomo II un bloque cuyo contenido principal consista en:
\begin{enumerate}
  \item refundar primitivas;
  \item repetir cláusulas basales sin ganancia operativa;
  \item introducir laboratorio no normalizado;
  \item depender de retórica donde falta estructura;
  \item anticipar claims de Tomo III sin contrato suficiente.
\end{enumerate}
\end{definition}

\begin{observation}
Esta condición de exclusión protege la función intermedia del tomo y evita que el canon operativo se vuelva un depósito heterogéneo.
\end{observation}

\section{Ganancia operativa}
\begin{definition}
Se entiende por ganancia operativa todo aumento explícito de capacidad para:
\begin{enumerate}
  \item distinguir niveles,
  \item fijar umbrales,
  \item ordenar transiciones,
  \item comparar configuraciones,
  \item preparar protocolos o criterios utilizables por bloques posteriores.
\end{enumerate}
\end{definition}

\begin{observation}
La ganancia operativa no equivale a incremento verbal. Debe expresarse como una mejora real de tratamiento estructural o jerárquico.
\end{observation}

\section{Cláusula de herencia}
\begin{observation}
Toda pieza admitida en el Tomo II debe declarar, al menos de forma implícita pero legible, de qué parte del núcleo basal integrado en Tomo I hereda su sentido estructural.
\end{observation}

\begin{observation}
Cuando una pieza no pueda expresar esa herencia, debe revisarse su ubicación o diferirse hasta contar con dependencia canónica visible.
\end{observation}

\section{Cláusula de alcance}
\begin{observation}
El contrato operativo fija el umbral editorial y formal de admisión al canon intermedio. No sustituye las pruebas específicas ni las validaciones particulares que puedan exigirse a bloques posteriores.
\end{observation}



